(1) 编码态同时满足 \(S_1=S_2=+1\)，因为 \(|000\rangle,|111\rangle\) 的相邻 \(Z\) 本征值乘积均为正。利用 \(ZX=-XZ\)，综合征依次为：\(I:(+,+)\)，\(X_1:(-,+)\)，\(X_2:(-,-)\)，\(X_3:(+,-)\)。测量仅提取错误子空间标签，对任意 \(\alpha,\beta\) 都是确定综合征，不测量逻辑振幅。
(2) 零或一次翻转可由综合征唯一纠正；两次翻转与剩余那一位单错具有相同综合征，恢复后变成 \(X_1X_2X_3=X_L\)。三次翻转也直接为 \(X_L\)。因此
\[
P_L=\binom32p^2(1-p)+p^3=3p^2-2p^3.
\]
(3) 逻辑通道 \(\mathcal E(\rho)=(1-P_L)\rho+P_LX_L\rho X_L\)。对标准Bell态，\(X\) 将其映射到正交Bell态，故纠缠保真度 \(F_e=1-P_L\)。未编码保真度 \(1-p\)。差 \(p-P_L=p(1-p)(1-2p)\)，所以严格改善当且仅当 \(0<p<1/2\)。\(p=0,1/2,1\) 两者相等，极限均物理合理。
(4) 所有 \(Z_i\) 与 \(S_1,S_2\) 对易，测量结果仍 \((+,+)\)；在编码空间 \(Z_1=Z_2=Z_3=Z_L\)，偶数个 \(Z\) 相互抵消，奇数个留下逻辑相位翻转：
\[
P_Z=\binom31q(1-q)^2+q^3=3q-6q^2+4q^3.
\]
低噪声时 \(P_Z\simeq3q\) 比单比特更差。此码只保护一个Pauli噪声轴；同时防 \(X,Z\) 需要互补基编码或例如Shor、Steane等更完整的码。
